\section{Latihan}
\label{sec:latihan-bab6}

\begin{enumerate}
    \item \textbf{Konseptual.} Jelaskan mengapa 8086 yang berbus
          alamat 20 bit memakai register segmen 16 bit. Apa yang
          akan hilang bila prosesor hanya memakai satu register
          16 bit untuk menunjuk alamat?

    \item \textbf{Perhitungan.} Hitung alamat fisik dari pasangan
          berikut. Tuliskan langkahnya.
          \begin{enumerate}
              \item \code{0000h:1A00h}
              \item \code{1500h:0100h}
              \item \code{9FFFh:000Fh}
              \item \code{B800h:0060h}
              \item \code{ABCDh:1234h}
          \end{enumerate}

    \item \textbf{Overlapping.} Alamat fisik \hex{00080} dapat
          dinyatakan oleh beberapa pasangan \eng{segment:offset}.
          \begin{enumerate}
              \item Tentukan banyaknya pasangan dengan pola
                    $\lfloor F/16 \rfloor + 1$.
              \item Tuliskan seluruh pasangan tersebut.
              \item Jelaskan mengapa jumlah pasangan untuk
                    \hex{00080} lebih banyak daripada \hex{00050}.
          \end{enumerate}

    \item \textbf{Peta memori.} Berdasarkan
          Tabel~\ref{tab:peta-memori}, jawab pertanyaan berikut.
          \begin{enumerate}
              \item Wilayah mana yang disebut memori konvensional,
                    dan berapa ukurannya?
              \item Apa isi area \hex{A0000}--\hex{BFFFF}?
              \item Di wilayah mana ROM BIOS berada?
          \end{enumerate}

    \item \textbf{Penulisan program.} Tulislah program yang
          mendeklarasikan tiga variabel pada segmen data, lalu
          menampilkan offset ketiganya dengan operator
          \code{OFFSET}. Gunakan prosedur \code{cetak\_hex} pada
          Kode~\ref{lst:alamat-fisik} sebagai acuan.

    \item \textbf{Penulisan program.} Tulislah program yang
          menghitung alamat fisik pasangan \code{1234h:5678h}
          secara langsung, yaitu memuat kedua nilai sebagai
          konstanta ke register, lalu melakukan pergeseran dan
          penjumlahan seperti pada Kode~\ref{lst:hitung-alamat}.
          Program harus mencetak \hex{179B8}.

    \item \textbf{Analisis kesalahan.} Seorang mahasiswa menghitung
          alamat fisik \code{1234h:5678h} dan memperoleh
          \hex{68AC}. Ia juga melihat DEBUG menampilkan
          \code{076A:0000} dan menyimpulkan alamat fisiknya
          \hex{0000}.
          \begin{enumerate}
              \item Tunjukkan kesalahan pada perhitungan pertama
                    dan hitung hasil yang benar.
              \item Jelaskan kesalahan penalaran pada kesimpulan
                    kedua, lalu tentukan alamat fisik yang benar
                    bila \reg{DS} $= \hex{076A}$.
              \item Sebutkan satu cara memeriksa kebenaran hasil.
          \end{enumerate}

    \item \textbf{Konseptual.} Jelaskan hubungan antara 1 MB,
          segmen 64 KB, dan paragraf 16 byte. Mengapa satu alamat
          fisik dapat memiliki paling banyak 4.096 pasangan
          \eng{segment:offset}?
\end{enumerate}

\begin{refleksibox}
\vspace{1.1em}
Jawab pertanyaan berikut pada jurnal praktikum Anda.
\begin{enumerate}
    \item Bagian mana yang paling sulit: aritmetika heksadesimal,
          istilah segmentasi, atau pembacaan peta memori? Mengapa?
    \item Setelah membuktikan \eng{overlapping} dengan DEBUG,
          bagaimana pandangan Anda tentang hubungan alamat logis
          dan alamat fisik?
    \item Konsep apa yang perlu Anda ulangi sebelum mempelajari
          memulai assembly pada bab berikutnya?
\end{enumerate}
\end{refleksibox}
